\section{Future Work}\label{sec:conc}

Because \skye\ is a framework for developing new EOS physics we expect future work to bring several key improvements.
First, and most pressing, is handling of partial ionization and neutral matter.
With that \skye\ could be used across the entire range of densities and temperatures which arise in stellar evolution calculations.
This could be done in a Debye-Huckle-Thomas-Fermi formalism~\citep{PhysRev.111.1460} or other approaches in the physical picture~\citep{Rogers2002}, or else via free energy minimization~\citep{Irwin2004} in the chemical picture~\citep{Saumon1995}.
The key constraint in each of these approaches is that \skye\ needs to remain fast enough to use in practical stellar evolution calculations.
Our hope is that the flexibility afforded to \skye\ by its automatic differentiation machinery will allow us to rapidly prototype and test these various possibilities.

Along similar lines, \skye\ could be made to support phase separation by minimizing the free energy with respect to the compositions of the liquid and solid phases.
The major bottleneck to supporting this is the current lack of Fortran compiler support for parameterized derived types.
Once this compiler challenge is resolved, phase separation physics  should not be difficult to implement.

More broadly, we make \skye\ openly available with the hope that it will grow into a community resource
to use automatic differentiation to explore analytic free energy terms that captures improvements in
existing physics and development of new or not yet considered physics.
